\section{微分同胚}
\label{sec:mrel-diffeomorphisms}

两个流形什么时候应当看作“同一个”光滑流形？集合论只问元素能否一一对应，拓扑只问邻近关系能否一一对应，而光滑流形还带有光滑结构，即哪些函数光滑、哪些曲线有速度。要说两个流形在光滑意义下相同，就要求有一个对应把这些信息双向保留下来。

\paragraph{定义}
保留光滑结构，意味着对应本身光滑，并且反过来的对应也光滑。

\begin{definition}[微分同胚]
\label{def:mrel-diffeomorphism}
设$M,N$是光滑流形。若$F:M\to N$是光滑双射，且逆映射$F^{-1}:N\to M$也光滑，则称$F$为\term{微分同胚}（diffeomorphism）。存在微分同胚时，称$M$与$N$微分同胚。
\end{definition}

微分同胚$F$把$N$上的光滑函数$g$对应到$M$上的光滑函数$g\circ F$，反过来$f\mapsto f\circ F^{-1}$，两种对应互逆。所以在两个流形上能谈论的“光滑对象”完全相同，只是名字换了。第~\ref{chap:manifold-maps}章又说明，光滑双射且逆光滑的映射把向量场一一推送为光滑向量场，所以微分同胚也保留了向量场。

\paragraph{三种对应的比较}
集合双射、同胚与微分同胚依次保留越来越多的信息，可以用例子看清每一层多出了什么。

只是集合双射，几乎什么都不保留。$\mathbb R$与$\mathbb R^2$的元素可以一一对应，但这样的双射不可能是同胚，因为维数是拓扑不变量；这一维数不变性需要较深的拓扑工具，这里不给出证明。同胚保留连通性、紧性这类拓扑性质，但不保留光滑性。$F(x)=x^3$是$\mathbb R$到自身的同胚，也是光滑双射，但逆映射$x^{1/3}$在原点不可微，所以$F$不是微分同胚。在$F$下，光滑函数$g(y)=y$对应到$g\circ F^{-1}$即$x^{1/3}$，不再光滑。

微分同胚是比较两个光滑结构的正确工具，而不是比较两个图册的写法。第~\ref{chap:manifold-basics}章指出，$\mathbb R$上的图$\chi(x)=x^3$与恒等图不光滑相容，记带$\chi$所决定结构的实轴为$\mathbb R_\chi$。恒等映射$\mathbb R\to\mathbb R_\chi$是同胚，但不是微分同胚。可是$G(x)=x^{1/3}$是从标准$\mathbb R$到$\mathbb R_\chi$的微分同胚：坐标表示是$\chi\circ G(x)=x$，逆映射的坐标表示是$G^{-1}\circ\chi^{-1}(z)=(z^{1/3})^3=z$，都光滑。所以这两个光滑结构不同，得到的光滑流形却微分同胚。

\paragraph{微分同胚的微分可逆}
微分同胚在每点的微分都是线性同构，这是链式法则的直接推论，也说明微分同胚的流形维数相同。

\begin{theorem}[微分同胚的微分]
\label{thm:mrel-diffeo-differential}
若$F:M\to N$是微分同胚，则对每个$p\in M$，$\mathrm dF_p:T_pM\to T_{F(p)}N$是线性同构，其逆为$\mathrm d(F^{-1})_{F(p)}$。特别地，$\dim M=\dim N$。
\end{theorem}

\begin{proof}
对$F^{-1}\circ F=\mathrm{id}_M$与$F\circ F^{-1}=\mathrm{id}_N$分别用链式法则，得$\mathrm d(F^{-1})_{F(p)}\circ\mathrm dF_p=\mathrm{id}_{T_pM}$与$\mathrm dF_p\circ\mathrm d(F^{-1})_{F(p)}=\mathrm{id}_{T_{F(p)}N}$。两个有限维空间之间存在互逆的线性映射，维数相同。
\end{proof}

\paragraph{反函数定理}
反过来，若只知道$\mathrm dF_p$可逆，能推出什么？在欧氏空间中，局部反函数定理说：Jacobian矩阵在一点可逆的连续可微映射，在该点附近有连续可微的逆。这一定理及隐函数定理可参考\BookChapterRef{chap:implicit-inverse-functions}{卷二第19章“隐函数与反函数”}；光滑映射的逆映射也光滑，因为逆映射的Jacobian矩阵$[\bm J_G]^{-1}$是光滑函数的复合，逐阶求导即可。借助坐标图，这一结论搬到流形上。

\begin{theorem}[流形上的反函数定理]
\label{thm:mrel-inverse-function}
设$M,N$是无边界光滑流形，$F:M\to N$光滑，$p\in M$，且$\mathrm dF_p$可逆。则存在$p$的开邻域$U$与$F(p)$的开邻域$V$，使$F|_U:U\to V$是微分同胚。
\end{theorem}

\begin{proof}
取包含$p$与$F(p)$的图，使$F(U_1)\subseteq V_1$。坐标表示$\hat F$在$\varphi(p)$处的Jacobian矩阵就是$\mathrm dF_p$的矩阵，可逆，于是两侧维数相同。由欧氏反函数定理，存在$\varphi(p)$的开邻域$\widehat U$与$\hat F(\varphi(p))$的开邻域$\widehat V$，使$\hat F:\widehat U\to\widehat V$是光滑双射且逆光滑。令$U=\varphi^{-1}(\widehat U)$，$V=\psi^{-1}(\widehat V)$，则$F|_U=\psi^{-1}\circ\hat F\circ\varphi$是三个微分同胚的复合。
\end{proof}

$\varphi(p)$在$\varphi(U_1)$的内部是这里不可缺少的。在半空间的边界点处，坐标表示的值域可能被迫落在半空间中，欧氏反函数定理给出的开邻域不再适用，所以本章把讨论限制在无边界情形。

\paragraph{局部可逆与整体双射}
若$F$在每点都满足定理~\ref{thm:mrel-inverse-function}的条件，就称它为\term{局部微分同胚}（local diffeomorphism）。局部微分同胚未必是微分同胚，因为反函数定理只给出局部的逆，不涉及整体是否一一对应。

映射$E:\mathbb R\to S^1$，$E(t)=(\cos t,\sin t)$，在角度图中坐标表示的导数为$1$，所以是局部微分同胚；但$E(t)=E(t+2\pi)$，它不是单射。圆周上的倍角映射$(\cos\theta,\sin\theta)\mapsto(\cos2\theta,\sin2\theta)$在角度坐标中是$\theta\mapsto2\theta$（相差$2\pi$的整数倍），导数为$2$，也是局部微分同胚，并且是满射，可每点恰有两个原像。这两个例子说明，导数处处可逆既不推出单射，也与满射无关。

缺的正是整体的一一对应。若局部微分同胚$F$还是双射，它就是微分同胚：逆映射$F^{-1}$在每点$q$附近与反函数定理给出的局部逆一致，局部逆光滑，所以$F^{-1}$光滑。类似地，单射的局部微分同胚是到其像的微分同胚，并且像是开集，因为每点的像都含有一个开邻域。

\paragraph{光滑结构相同不等于几何相同}
微分同胚只关心光滑结构，不关心长度与距离。映射$x\mapsto\tan(\pi(x-\tfrac12))$是$(0,1)$到$\mathbb R$的微分同胚，导数$\pi/\cos^2(\pi(x-\tfrac12))$处处为正，逆映射$y\mapsto\tfrac12+\tfrac1\pi\arctan y$也光滑。前者长度为$1$，后者无界。再如$x\mapsto2x$是$\mathbb R$到自身的微分同胚，把所有距离加倍。所以微分同胚的两个流形可以有完全不同的度量性质。保持长度的映射需要在切空间上先有内积，这样的等距映射在第~\ref{chap:riemannian-geometry}章讨论。
